\section{All binary Boolean gates}
Let $C_2$ count unrestricted binary Boolean gates. Constants are free.
\begin{theorem}\label{th:fullbasis}
There is an absolute $c_2>0$ such that, for every integer
$0\le s\le (1+c_2)n$, the total decision function
$\mathbf1[C_2(f)\le s]$ has bounded-error quantum query complexity
and ordinary $1/3$-approximate degree both $\Theta(\sqrt N)$.
\end{theorem}
\begin{proof}
Apply essential-variable recovery and handle $m\le n-\log\log(n+2)$ by the preceding identification and verification algorithm. The census remains valid with the larger fixed gate alphabet. Suppose henceforth $m=n-O(\log\log n)$ and $H=2^m$. We distinguish two subfamilies in the unlimited offline computation.

First take the family with algebraic degree at most $m-\lceil8\log n\rceil$ over $\mathbb F_2$. Two distinct such tables differ in at least $n^8$ entries, by the minimum distance bound $2^{m-d}$ for Boolean polynomials of degree at most $d$. One proof chooses a degree-$d'$ monomial of maximum degree in their difference: every restriction fixing its complement has odd weight on those $d'$ axes, yielding at least $2^{m-d'}$ ones. The sampling and identification argument of the preceding section, run on this subfamily with discrepancy threshold $n^6$, therefore recovers an \emph{exact} candidate on member inputs. Its verification on the original oracle uses $O(\sqrt N)$ queries.

For the other subfamily, any size-$s$ circuit on the $m$ essential axes has at most
\[
                        z\le s-m+1+\lceil8\log n\rceil
\]
affine gates after unary simplification. Indeed a function of algebraic degree $d$ needs at least $d-1$ nonlinear gates. For a direct proof replace every nonlinear gate by one AND with affine preprocessing and postprocessing. In a circuit with $a$ ANDs, the first AND multiplies two affine forms. If necessary make one nonconstant form an input by an invertible affine change of coordinates; restricting that input to either value eliminates the first AND. Induction gives degree at most $a$ on both restrictions, hence at most $a+1$ before restriction. Constant factors can simply be removed. An invertible affine change preserves degree. This is the standard multiplicative-complexity degree bound; see \href{https://pmc.ncbi.nlm.nih.gov/articles/PMC7047778/}{Boyar et al.}

Represent every essential nonlinear binary gate as an AND with possible complements on both inputs and on its output. Permit freely signed wires in this representation. Simplify gates depending on at most one input; each remaining affine gate is a signed XOR. In the trimmed graph the total excess fanout is at most $e=s-m+1$. The two structural lemmas used for the preceding reconstruction have the following variants.

\emph{Critical paths.} Stop each raw path at its first fork, output, or input port of an XOR, whichever comes first. There are at most $e+1+2z$ possible endpoints. Along a sensitive path before its endpoint, an enabled signed AND transmits its path bit with a sign fixed by the port and gate signs. Two distinct varying, simultaneously sensitive raw variables cannot arrive at the same such endpoint: at their first nonlinear merger both signed port bits would be forced to their enabling values on all witnesses. A first merger at an XOR uses two distinct stopping ports. Thus the two-witness path bound has right-hand side $e+1+2z$.

\emph{Sensitivity tail.} Cut each fork and each XOR. For $F$ the set of cut vertices, $|F|\le e+z$, the total incoming edges into $F$ number at most $2|F|$, and the total outgoing edges from $F$ number at most $|F|+e$. Excluding the at most $e$ raw fork variables, the remaining private raw variables lie in one-exit trees exactly as in Lemma~\ref{lem:tail}. For fixed incoming port values each signed read-once AND tree is unate in its private variables: changing the signs of the individual private inputs according to their unique paths turns it into a monotone function. The mean sensitivity bound $\sqrt{b_i}$ therefore holds for every fixed port setting. For explicit estimates put $e'=2(e+z+1)$. The previous bounds on the number $P$ of pieces and the total number of their entering ports hold with $4e'$ in place of $4e$ by the preceding incident-edge counts, and raw forks contribute at most $e'$. The rest of the independent-piece Hoeffding argument is unchanged. In particular the same exponential tail holds if $e'\le\gamma m$ for the absolute $\gamma$ of that proof.

Take, for example, $c_2=\gamma/16$. In the remaining regime $m=n-O(\log\log n)$, for sufficiently large $n$ both $e'\le\gamma m$ and $e+1+2z+m/8<m/3$. Run the reconstruction using only the high-degree subfamily for its initial representative and all projected candidate lists. The exponential tail controls the high-sensitivity set of the representative. The path bound for the unknown high-degree member puts all good discrepancy anchors in the same radius-$m/3$ ball. The two verifications, one for each candidate subfamily, accept a true member with bounded error and reject every nonmember with bounded error by checking candidate size and actual equality. Amplify each branch by a fixed constant if needed. This proves the upper bound; a quantum acceptance polynomial gives the approximate-degree upper bound.

For the lower bounds, duplicate one bit on each antipodal pair. The zero word has zero circuit cost, whereas a single marked pair has full-binary cost $2n-3$ by \href{https://doi.org/10.1016/0304-3975(76)90073-6}{Schnorr's equivalence theorem}, with arbitrary input signs absorbed into the first gates. Thus zero and every singleton give opposite labels for all $s\le 2n-4$. For completeness let $M=N/2$ and compose any degree-$d$ total approximator with this duplication. Its expectation $p(t)$ on independent Bernoulli-$t$ coordinates has degree at most $d$ and is bounded in absolute value by $4/3$ on $[0,1]$. At $t=a/M$ for a sufficiently small fixed $a>0$, the probability of one mark is $a(1-O(a))$ and the probability of at least two is $O(a^2)$. Every singleton differs in the same direction from the zero value by at least $1/3$. Hence $|p(a/M)-p(0)|=\Omega(1)$. The mean-value theorem and the ordinary Markov polynomial derivative bound on $[0,1]$ give $d^2=\Omega(M)$. The polynomial method then gives the same quantum lower bound.
\end{proof}

\begin{theorem}\label{th:fullconfidence}
With the same $c_2$, for $0<\delta\le1/3$ the full-binary total decision problem has bounded-error quantum query complexity and ordinary $\delta$-approximate degree
\[
  \Theta\left(\min\{N,\sqrt{N\log(1/\delta)}\}\right)
\]
uniformly for $0\le s\le(1+c_2)n$. The same bounds hold with counted De Morgan gates.
\end{theorem}
\begin{proof}
\setlength{\emergencystretch}{1em}
Write $\ell=\log(1/\delta)$. For $\ell\ge (\log n)/4$ use the high-confidence finite-family oracle identification bound from the companion manuscript \emph{The approximate degree of Boolean minimum circuit size}, section on high-confidence identification. Applied to the size-$s$ family, whose log cardinality is $O(n\log n)$, followed by candidate verification, it gives $O(\sqrt{N(\log n+\ell)})$. Alternatively one can read all $N$ bits. We give the remaining argument for $\ell<(\log n)/4$, where the desired failure probability is greater than $n^{-1/4}$.

Run the degree split directly on \emph{all} $n$ raw axes, without an essential-variable discovery stage. The low-degree class has distance at least $n^8$ on the full cube. In the high-degree class the number $m$ of essential inputs automatically lies between $n-\lceil8\log n\rceil$ and $n$. Its excess fanout and its number of XOR gates obey the same bounds as before with an additional $O(\log n)$ term. Each high-degree candidate has an exponentially small full-cube fraction of points of sensitivity above $n/32$: lifting its essential-axis tail does not change sensitivity. The two-anchor argument on the full cube remains valid; every raw variable counted as sensitive is necessarily essential. Thus the ball, the high-sensitivity region, and the punctured candidate-list search can all use the full truth-table addresses. This construction works for counted De Morgan candidates as well: maintain their exact candidate family offline and use their signed full-binary representations only for the structural estimates.

The only combinatorial failure not controlled by a query subroutine is selecting a remaining error with more than $n/32$ internal neighbors. Its probability is $O((\log n)/n)=o(\delta)$ because the search returns a uniform marked point conditioned on success. All identification stages on exponentially small regions cost $o(\sqrt N)$ even after $O(\log n)$ repetitions. The punctured list length is $D=O(N\log|\mathcal A|/n^6)=O(N/n^4)$ including a fixed enlargement of the sampling constant. The probability that a distant candidate agrees on the list can be made at most $\delta/100$ by using $D=\lceil 4N(\ln(|\mathcal A|+2)+\ell)/n^6\rceil$. Identification there costs at most $O(\sqrt{D\log(|\mathcal A|+2)})$, still $o(\sqrt N)$ after logarithmically many repetitions. This includes the low-degree candidate discovery. Use independent repetitions and plurality on the identified projections; on promised subroutine inputs the correct projection is unique. Offline index recovery and size checks remain unrestricted.

Only a constant number of searches on the entire cube remain, including final verification. A search with failure probability $O(\delta)$ costs $O(\sqrt{N(1+\ell)})$ queries and returns a uniform marked address conditioned on success. To see this, run exact amplitude amplification for each candidate mark count $j=1,\ldots,\lceil C(1+\ell)\rceil$, verifying every proposed address, at total cost $O(\sum_j\sqrt{N/j})$. For larger counts use $O(1+\ell)$ independent bounded-cost search attempts calibrated at the lower bound $C(1+\ell)$, allowing a uniform sample to handle a high marked fraction. Permutation symmetry gives the uniformity assertion. Summing the error of the constant number of global calls and amplified local stages proves the upper bound.

For the matching lower bound let $M$ count the easy tables. While at least two remain, fix a nonconstant coordinate to its rarer value among the survivors. After at most $\log M=O(n\log n)$ steps exactly one easy table remains. On the free coordinates the decision function is therefore equality to one specified word, equivalent by input and output signs to OR on $N-O(n\log n)$ bits. Its $\delta$-approximate degree is $\Omega(\min\{N,\sqrt{N\log(1/\delta)}\})$ by the small-error OR bound (see \href{https://arxiv.org/abs/0802.1816}{de Wolf, 2008}), including smaller errors by monotonicity. The polynomial method transfers this lower bound to quantum queries and the proved quantum upper bound to approximate degree.
\end{proof}
